% =====================================================================
%  Literature review: Teacher/school recommendations, asymmetric
%  information, and student choice.
%  Compiled as a reading list. Each \entry has a commented \note{}
%  slot right after it -- uncomment and fill in as you re-read.
%  Compile with: pdflatex litreview.tex
% =====================================================================
\documentclass[11pt]{article}
\usepackage[a4paper,margin=2.1cm]{geometry}
\usepackage[T1]{fontenc}
\usepackage[utf8]{inputenc}
\usepackage{xcolor}
\usepackage{titlesec}
\usepackage{parskip}
\usepackage{enumitem}
\usepackage[hidelinks]{hyperref}

% ---- colors ----
\definecolor{dark}{HTML}{1A2A40}
\definecolor{accent}{HTML}{3A5A78}
\definecolor{tagcol}{HTML}{3A5A78}
\definecolor{notecol}{HTML}{8A5A00}
\definecolor{grey}{HTML}{5B6770}

% ---- section formatting ----
\titleformat{\section}{\large\bfseries\color{dark}}{}{0pt}{}
\titlespacing*{\section}{0pt}{14pt}{2pt}
\titleformat{\subsection}{\normalsize\bfseries\color{accent}}{}{0pt}{}

% ---- helper macros ----
% \entry{citation}{summary}{TAG}
\newcommand{\entry}[3]{%
  \par\medskip\noindent\textbf{#1}\par
  \nopagebreak #2\hspace{0.35em}{\small\color{tagcol}[#3]}\par}
% annotation slot: uncomment and fill
\newcommand{\note}[1]{\par\smallskip\noindent\hangindent=1.3em
  \textcolor{notecol}{\footnotesize$\hookrightarrow$~\textit{#1}}}
% short section intro
\newcommand{\secintro}[1]{\par\smallskip\noindent\textcolor{grey}{\small #1}\par
  \vspace{2pt}\hrule height0.4pt\vspace{4pt}}

\setlength{\parindent}{0pt}

\begin{document}

% =====================================================================
\begin{center}
{\LARGE\bfseries\color{dark} Teacher and School Recommendations,\\[2pt]
Asymmetric Information, and Student Choice}\\[6pt]
{\color{grey}\itshape A panoramic, annotated reading list --- emphasis on empirical /
reduced-form work, with key structural and theory anchors flagged.}\\[2pt]
{\color{grey}\small Project on track recommendations in the Belgian (LOSO) setting.}
\end{center}
\vspace{4pt}
\hrule
\vspace{8pt}

\noindent\textbf{How to read / annotate this.} Organised around four questions:
(A) do schools/teachers hold information students and parents lack;
(B) why recommendations should be useful (causal / predictive / matching value);
(C) how schools form recommendations (determinants, bias, statistical discrimination);
(D) how students/parents respond. Sections E--I add the class-rank channel, the
belief-formation framework, the statistical-discrimination toolkit, tracking
foundations, and reviews, for context. \emph{Tags:}
\textbf{RF} reduced-form empirical; \textbf{EXP} field/lab/survey experiment;
\textbf{STR} structural; \textbf{THY} theory; \textbf{REV} review.
After each entry there is a commented \verb|% \note{}| line in the source ---
uncomment it and type your annotation; it renders as a brown indented note.
A few venues are flagged \emph{(verify)} where the published outlet/year should be
double-checked before citing.

% =====================================================================
\section{A. Do schools and teachers know more than students and parents?}
\secintro{Evidence that subjective teacher assessments carry information beyond
test scores, and that parents/students hold inaccurate beliefs an institution
could (or does) move.}





\entry{Jackson, C. K. (2018). What Do Test Scores Miss? The Importance of Teacher
Effects on Non--Test Score Outcomes. \textit{Journal of Political Economy}
126(5).}
{Teacher effects on behaviours (absences, suspensions, course grades, grade
progression) are only weakly correlated with their effects on test scores, yet
predict high-school completion and other long-run outcomes more strongly. Core
evidence that what teachers observe/affect beyond scores is outcome-relevant ---
the premise behind a school ``information set''.}{RF}
\note{This paper just is Teacher Value Added but to noncognitive outcomes. It shows that some teachers impact on cognitive scores is not that correlated with the impact on non-cognitive ---here noncognitive are measured as absentism, nonshowing up, which are correlated with grit and big5--- and that non-cognitive tend to be predictive of future outcomes even conditioning on test scores.}

\entry{Kinsler, J., and Pavan, R. (2016). Parental Beliefs and Investment in
Children: The Distortionary Impact of Schools. HCEO Working Paper 2016-029,
Univ.\ of Chicago.}
{Parents' beliefs about their child's skill relative to same-age peers are
distorted by the child's rank \emph{within their own school} (big-fish-little-pond):
parents in high-achieving schools underestimate, those in low-achieving schools
overestimate, and this feeds investment (homework help, tutoring). The
``Kinsler/Pavan'' reference --- schools shaping the information households act on.}{RF / STR}
 \note{this is a key reference. Here the information assumption is that parents don't have perfect information on the relative standing of the kids skills within the overall population and that this affects their investment. Parents know the relative within school and receive grades and teacher reports as signals that in part reflect the overall relative standing, update their priors etc and readjust investment. Parents of high SES/skill school overinvest and parents of low-skill school who are at the top underinvest\dots}

\entry{Dizon-Ross, R. (2019). Parents' Beliefs about Their Children's Academic
Ability: Implications for Educational Investments. \textit{American Economic
Review} 109(8): 2728--65.}
{Field experiment in Malawi: parents' baseline beliefs are often wrong and
information frictions are worse for the poor. Clear, digestible performance
information shifts beliefs and reallocates investment toward better-matched inputs.
Clean causal evidence that an institution holding accurate performance information
changes household decisions.}{EXP}
 \note{same line as in Bergman, parents hold inaccurate beliefs about performance and information interventions can have positive impact on adjustments in terms of effort etc.}

\entry{Boneva, T., and Rauh, C. (2018). Parental Beliefs about Returns to
Educational Investments --- The Later the Better? \textit{Journal of the European
Economic Association} 16(6): 1669--1711.}
{Elicits parents' beliefs about the \emph{technology} mapping investments to child
outcomes; beliefs vary widely, predict actual investment, and differ by SES (lower-SES
parents perceive lower returns). Directly supports modelling SES-graded priors over
the skill-to-outcome mapping.}{EXP / survey}
\note{Again, focused in explaining differential parental investments on child based on differences in beliefs. This only informs that priors should differ by SES or origin of the family. Main finding is that beliefs of parents go against evidence that early childhood investments are more important than late and that there are complementarities. }


\entry{Eva Feron, Trudie Schils e.g.\ ``The predictive power of track
recommendations in Dutch secondary education'' and ``Does the teacher beat the test?
The value of the teacher's assessment in predicting student ability''.}
{Dutch micro-data finding the subjective teacher assessment is roughly twice as
predictive as primary-school cognitive scores for initial track placement, and more
predictive of later allocation than scores. Strong ``teachers know more'' evidence,
but confounds prediction with self-fulfilment. Confirm authors/venue.}{RF (verify)}
 \note{important reduced form paper that, in a very similar context as in Belgium, shows that recommendations are more predictive than objective grades for future success ---also points to schools observing daily the students and incorporating noncognitive assessment in top of their recommendations (i.e grit, perseverance etc, which are fundamental skills for certain tracks). Importantly, it points to the subtle issue of counfounding predictiveness with self-fulfillment: recommendation can be predictive or it can also change the behavior of the kid (Pygmalion theory that students react to expectations of the teacher)}

\entry{[Bruno Ferman, Luiz Felipe Fontes (2022)] Assessing Knowledge or Classroom
Behaviour? Evidence of Teachers' Grading Bias (Journal of Public Economics).}
{Exploits staggered introduction of conduct grades across German states to estimate
causal effects on school-to-work transitions and non-cognitive skills, and documents
that subject grades themselves encode behaviour. Speaks to the soft (non-test) content
teachers fold into evaluations. Verify authorship.}{RF (verify)}
\note{We observe objective grades, and then assessment of teachers on performance etc. Again, it supports the idea that teachers evaluate student's fit for the different track based on their noncognitive skills on top of cognitive skills. They write something like his argument fits naturally within a principal-agent framework of strategic information transmission, where transmitting biased information leads agents to take suboptimal decisions. Indeed, Mechtenberg (2009) shows in a cheap talk model that distorting grades based on student behavior may discourage students from making human capital investments depending on how they interpret the meaning of their grades. }

\entry{Bergman (2021), "Parent-Child Information Frictions," JPE}
{.}{RF }
\note{parents have imperfect information of the effort of the student in class and o }



Afegir:
Bergman (2021), "Parent-Child Information Frictions," JPE 
I think we can think of this as effort in a given task. 
% =====================================================================
\section{B. Why should recommendations be useful? Causal and matching value}
\secintro{Whether teacher expectations/recommendations causally improve (or distort)
outcomes, whether discretion adds value as a matching device, and how a teacher mark
signals within the system.}

\entry{Papageorge, N. W., Gershenson, S., and Kang, K. M. (2020). Teacher
Expectations Matter. \textit{Review of Economics and Statistics} 102(2): 234--251.}
{Two teachers reported expectations for the same student; their disagreement (a
measurement-error design) identifies a causal effect. Teacher expectations raise the
probability of college completion (elasticity $\approx 0.12$); teachers are on average
over-optimistic, with biases acting like self-fulfilling prophecies. The key
``expectations are causal'' paper.}{RF}
\note{Issue is that expectations are a function of $\theta_i$ which also affects students outcomes. Identification of causal effects of teacher expectations (I think x will graduate inX on the prob) hinges on observing 2 diverging assessment for the same student. Two mechanisms, are they explores? So it is self-fullfiling or correct prediction. They want to answer simply how biased are teachers with respect to the truth\dots teacher bias arises when a teacher’s expectations diverge from information that is common to both of them, including ninth-grade GPA and the latent factor $\theta_i$. allow teachers to be wrong on average and allow them to be wrong on how they map $\theta_{i}$ to outcomes. i.e they find teachers tend to overestimate the chances and be overoptimisitc. For them, the main goal of the paper is to recover the distribution of teacher biases. }

\entry{Gershenson, S., Holt, S. B., and Papageorge, N. W. (2016). Who Believes in
Me? The Effect of Student--Teacher Demographic Match on Teacher Expectations.
\textit{Economics of Education Review} 52.}
{Two teachers' differing expectations for the same student reveal systematic
demographic-match effects (e.g.\ non-Black teachers hold lower expectations for Black
students). Underpins the idea that recommendations embed the evaluator's identity/
priors, not just the student's signal.}{RF}
\note{reduced form findings that exploit same argument as previous paper.}

\entry{Diamond, R., and Persson, P. (2016). The Long-term Consequences of Teacher
Discretion in Grading of High-stakes Tests. NBER Working Paper 22207.}
{Bunching in Swedish math-test distributions shows teachers ``bump up'' students who
had a bad test day (no gender/immigrant discrimination here). A Wald estimator shows a
higher discretionary grade raises later grades, graduation, college entry, and earnings
--- even though grades don't raise human capital directly. Interprets grades as a signal
\emph{within} the education system, with effort/perceived-ability complementarities.
Close to your recommendation-as-signal mechanism.}{RF}
note{natural question is whether school discretion to act beyond objective measures is corrective --student should have passed-- or increases biases etc Completely off from our interests I think}

\entry{Mulhern, C. (2022). Beyond Teachers: Estimating Individual Guidance
Counselors' Effects on Educational Attainment. EdWorkingPaper 22-632 (Annenberg/Brown).}
{Uses quasi-random assignment of students to counsellors in Massachusetts; counsellors
vary substantially in their effect on graduation, college attendance, selectivity, and
persistence. Direct evidence that an advisor holding information/guidance moves
attainment. Verify published version.}{RF (verify)}
% \note{}

\entry{Ichino, A., Mealli, F., and Viviens, J. (2025). Do Test Scores Help Teachers
Give Better Track Advice to Students? A Principal Stratification Analysis. Working
paper (arXiv 2511.05128 / CRENoS).}
{Dutch quasi-randomised cutoff lets teachers revise recommendations upward when scores
beat expectations; allowing upgrades raises successful placement in demanding tracks by
$\geq$6\% while misplacing $\approx$7\% of weaker students, with a welfare analysis over
short- vs long-term losses. Directly on the value of combining discretion with objective
signals.}{RF}
\note{The closest existing answer to "what is a good recommendation?", using a Dutch quasi-experiment (teachers may upgrade when the central test clears a cutoff) and a principal-stratification metric. It matters for three reasons: it supplies a welfare-relevant definition of recommendation quality (direct toward the hardest track a student can complete; minimize costly track changes); its fairness results — test-score information induces larger upgrades for immigrant and low-SES students, interpreted as "a larger shift in teachers' posterior beliefs" — are precisely the statistical-discrimination-through-beliefs and preference-contamination channels we model, observed in the field; and its robustness section treats the direct effect of the signal on the student (confidence, parental advocacy) as a nuisance to be bounded away, whereas that channel is the entire object of our model. It is also the sharpest competitor to position against: it is reduced-form, single-regime, and has no decision-making student, so it cannot run counterfactuals on the information structure or the school's objective — which is where our structural model contributes..}


\entry{Simon Burgess and Ellen Greaves, Test Scores, Subjective Assessment and Stereotyping  of Ethnic Minorities, JOLE (2013)}{We assess whether ethnic minority pupils are subject to low teacher expectations. We exploit the English testing system of “quasi-blind” externally marked tests and “non-blind” internal assessment to compare differences in these assessment methods between White and ethnic minority pupils. We find evidence that some ethnic groups are systematically “under-assessed” relative to their White peers, while some are “over-assessed”. We propose a stereotype model in which a teacher’s local experience of an ethnic group affects assessment of current pupils; this is supported by the data}{rt}
\note{When forming an assessment of a pupil’s likely progress, teachers use information on the past performance of members of that group in that school from previous years. The dependence of a pupil’s assessment on the performance of others of her ethnic group locally means that school composition matters. This is a form of indirect peer effect, and suggests another basis for parents and pupils selecting particular schools. Marginal reference, it deals with assessment as in grades and is centered around statistical discrimination. Model includes or accounts for past perceentage of people of the same background.}


\entry{Aderonke Osikominu, Gregor Pfeifer, Kristina Strohmaier, The Effects of Free Secondary School Track Choice: A Disaggregated Synthetic Control Approach, 2021 }{We exploit a recent state-level reform in Germany that granted parents the right to decide on the highest secondary school track suitable for their child, changing the purpose of the primary teacher's recommendation from mandatory to informational. Applying a disaggregated synthetic control approach to administrative district-level data, we find that transition rates to the higher school tracks increased substantially, with stronger responses among children from richer districts. Simultaneously, grade repetition in the first grades of secondary school increased dramatically, suggesting that parents choose school tracks also to align with their own aspirations – resulting in greater misallocation of students.}{rt}
\note{The cleanest causal evidence that the binding-versus-informational status of a recommendation has large allocative consequences: a German reform that made recommendations non-binding raised transitions to higher tracks (concentrated among high-SES districts) and nearly doubled grade-five repetition, read as aspiration-driven misallocation. It matters because it empirically validates the authority channel we otherwise risk treating as a footnote, and shows its direction — teacher authority was correcting an aspiration-driven distortion, not creating one. Crucially, it infers this mechanism indirectly, from the SES gradient plus the repetition jump; it observes neither beliefs nor the teacher's information, so it cannot separate "parents are better informed" from "parents chase aspirations." That separation — via our belief panel and terminal outcomes — is our contribution, in a setting where a soft advisory recommendation and a binding certificate operate jointly on the same student rather than across a reform, and where the recommending body itself bears the downstream repetition cost.}

Same study of the reform in Maximilian Bach's Heterogeneous responses to school track choice: Evidence from the repeal of binding track recommendations


\entry{Anne van Leest , Lisette Hornstra, Jan van Tartwijk and Janneke van de Pol, Test- or judgement-based school track recommendations: Equal opportunities for students with different socio-economic backgrounds?, 2021 }{There are concerns that school track recommendations that are mostly based on teachers’ judgements of students’ performance (‘judgement-based recommendations’) are more biased by students’ SES than school track recommendations that are mostly based on standardized test results (‘test-based recommendations’). A recent policy reform of the Dutch educational system has provided us the unique opportunity to compare the effects of students’ SES on these two types of track recommendations.  Aims. The aim of this study was to examine the differences between test-based and judgement-based recommendations regarding the direct and indirect effect of students’ SES at student level and school level.}{rt}
\note{There are valid arguments in favour of both types of recommendations. A central argument in favour of using test-based recommendations is that using a school leavers’ test as primary indicator improves educational equality. A school leavers’ test consists of the same set of questions and is ideally administered under similar conditions (Knoester Au, 2017). Additionally, for all students, the standard is set at the same level to ensure that they are judged solely on their performance, neglecting other student characteristics. Consequently, test-based recommendations should be similar for all students with comparable performance levels, regardless of their background characteristics. In turn, students from different backgrounds will have equal opportunities to be assigned a certain track recommendation (OECD, 2016). However, a school leavers’ test is administered at one specific moment, which makes it impossible to take students’ (cognitive) development throughout primary school into account (Driessen, 2005; OECD, 2016). When a student does not perform as well as he or she normally does, the results of the school leavers’ test are not in line with the students’ actual abilities. A central argument in favour of using judgement-based recommendations is that teachers have the opportunity, in addition to standardized test results, to include (non-cognitive) information that may be predictive of students’ future secondary school success, such as classroom behaviour, motivation, talents, development, work attitude, and school engagement. Consequently, in educational systems that use judgement-based recommendations, students who perform at the same educational level may receive different track recommendations. Moreover, teacher judgements are susceptible to bias, which may cause an undesired effect: lower track recommendations for students with more disadvantaged backgrounds.
}


\entry{De Ree, Oosterveen, Webbink, The quality of school track assignment decisions by  teachers, 2025 }{This paper analyzes the effects of educational tracking and the quality of track assignment decisions. We motivate our analysis using a model of optimal track assignment under uncertainty. This model generates predictions about the average effects of tracking at the margin of the assignment process. In addition, we recognize that the average effects do not measure noise in the assignment process, as they may reflect a mix of both positive and negative tracking effects. To test these ideas, we develop a flexible causal approach that separates, organizes, and partially identifies tracking effects of any sign or form. We apply this approach in the context of a regression discontinuity design in the Netherlands, where teachers issue track recommendations that may be revised based on test score cutoffs, and where in some cases parents can overrule this recommendation. Our results indicate substantial tracking effects: between 40\% and 100\% of reassigned students are positively or negatively affected by enrolling in a higher track. Most tracking effects are positive, however, with students benefiting from being placed in a higher, more demanding track. While based on the current analysis we cannot reject the hypothesis that teacher assignments are unbiased, this result seems only consistent with a significant degree of noise. We discuss that parental decisions, whether to follow or deviate from teacher recommendations, may help reducing this noise.}{rt}
\note{just shows how marginal students at the cutoff who go to a better track have better outcomes. Pretty obvious. }


If recommendations carry indeed ordinal information that students ignore:

\entry{Ahmat, Iriberri, The importance of relative performance feedback information: Evidence from a natural experiment using high school students, 2010, Journal of Public Econ }{We study the effect of providing relative performance feedback information on performance, when individuals are rewarded according to their absolute performance. A natural experiment that took place in a high school offers an unusual opportunity to test this effect in a real-effort setting. For one year only, students received information that allowed them to know whether they were performing above (below) the class average as well as the distance from this average. We exploit a rich panel dataset and find that the provision of this information led to an increase of 5\% in students' grades. Moreover, the effect was significant for the whole distribution. However, once the information was removed, the effect disappeared. To rule out the concern that the effect may be artificially driven by teachers within the school, we verify our results using national level exams (externally graded) for the same students, and the effect remains.
}{rt}
\note{A more detailed analysis shows that there are heterogeneous effects from the provision of relative performance feedback information. First, the effect is significant only for students in the first and fourth years of high school. It is reasonable to assume that in the first year of high school the relative performance feedback information provides new information and that students might be more reactive to it than in the subsequent years. As for the fourth year effect, given that the grades during this year are especially important in determining the final university entry grade, one may believe that any additional information regarding grades will provoke a stronger reaction. Second, the positive effect is strongest in science subjects such as Mathematics, as well as in language subjects. This has important policy implications since there has been a special interest in improving grades in more technical subjects such as Math. Third, we also find that, although girls overall obtain better grades than boys there is no significant gender difference in the reaction to the relative performance feedback information.9
 }


 \entry{Elsner, Isphording, big fish in a small pond:ability rank and human capital investment}{We study the impact of a student's ordinal rank in a high school cohort on educational attainment several years later. To identify a causal effect, we compare multiple cohorts within the same school, exploiting idiosyncratic variation in cohort composition. We find that a student's ordinal rank signiffcantly affects educational outcomes later in life. If two students with the same ability have a different rank in their respective cohort, the higherranked student is signiffcantly more likely to finish high school, attend college, and complete a 4-year college degree. These results suggest that low-ranked students under-invest in their human capital even though they have a high ability compared to most students of the same age. Exploring potential channels, we end that students with a higher rank have higher expectations about their future career, and feel that they are being treated more fairly by their teachers.
}{rt}
\note{ In theory, this result can be explained by at least four mechanisms. First, the rank may provide students with a noisy signal of their own ability. Students may conclude from a low relative ability that they have a low absolute ability. If a low perceived ability translates into low expected returns to college, students may choose not to go to college. Second, rank can affect intrinsic factors. Students with a higher rank may be more motivated and self-confident, and hence put more effort into their studies, which then translates into a higher chance of going to college. Third, a student's environment may be responsive to a student's rank. Teachers, family, and friends may offer more support to high-ranked students, leading to better grades and a higher chance of going to college. Finally, the result could be explained by selective college admission policies. Colleges often observe a student's GPA rank within a cohort, which is correlated with our ability measure. If admissions officers give priority to students with a higher rank regardless of the school quality, or if colleges automatically admit the top 10\% of a school cohort, then this can explain the effect.
 }

  \entry{Richard Murphy and Felix Weinhardt, Top of the Class: The Importance of Ordinal Rank}{This paper establishes a new fact about educational production: ordinal academic rank during
primary school has long-run impacts that are independent from underlying ability. Using data on
the universe of English school students, we exploit naturally occurring differences in achievement
distributions across primary school classes to estimate the impact of class rank conditional on
relative achievement. We find large effects on test scores, confidence and subject choice during
secondary school, where students have a new set of peers and teachers who are unaware of the
students’ prior ranking. The effects are especially large for boys, contributing to an observed
gender gap in end-of-high school STEM subject choices. Using a basic model of student effort
allocation across subjects, we derive and test a hypothesis to distinguish between learning and
non-cognitive skills mechanisms and find support for the latter.
}{rt}
\note{ This higher confidence could be indicative of two mechanisms. First, confidence could be reflective of students learning about their own strengths and weaknesses in subjects, similar to Azmat and Iriberri (2010) or Ertac (2006) where students use their test scores and cardinal relative
positions to update beliefs. Alternatively, consistent with the Big Fish Little Pond effect, which
has been found in many countries and institutional settings (Marsh et al., 2008), confidence due to
rank improves non-cognitive skills and lowers the cost of effort in that subject. The idea is that,
when surrounded by people who perform a task worse than oneself, one develops confidence in
that area. So, if a student is confident in her maths skills, she will be more resilient and have more
grit in solving maths exercises compared to another student of the same ability but different confidence. Using a stylized model of students trying to maximise test scores for a given total effort
and ability level across subjects, we derive a test to distinguish the Big Fish Little Pond from the
leaning mechanism and find evidence in favor of the former.
 }



% =====================================================================
\section{C. How do schools form recommendations? Determinants, bias, discrimination}
\secintro{The empirical core on what drives recommendations/grades and whether they are
biased by gender, ethnicity, or SES --- including blind-vs-non-blind designs built to
isolate bias.}

\entry{Carlana, M. (2019). Implicit Stereotypes: Evidence from Teachers' Gender
Bias. \textit{Quarterly Journal of Economics} 134(3): 1163--1224.}
{Exploiting near-random assignment of students to teachers with different Gender-Science
IAT scores: teacher stereotypes widen the math gender gap and push girls toward less
demanding high schools, \emph{following the teacher's track recommendation}, partly via
lowered female self-confidence in math. The model paper for bias working through
recommendations and beliefs.}{RF}
\note{also seen in the data, but hard to separate in our case from prefrences of girls, as it is data also from 1990s and gender norms are likely much more salient.}

\entry{Carlana, M., La Ferrara, E., and Pinotti, P. (2022). Implicit Stereotypes in
Teachers' Track Recommendations. \textit{AEA Papers and Proceedings} 112: 409--414.}
{Track-recommendation-specific companion to Carlana (2019): teacher implicit bias maps
into the recommendations minority/immigrant students receive.}{RF}
\note{ the families most subject to biased signals are the ones most reliant on them. Teachers’ expectations play an important role in shaping students’ educational
careers and may lead to a self-fulfilling
prophecy if children internalize low beliefs
on their educational success}


\entry{Lavy, V. (2008). Do Gender Stereotypes Reduce Girls' or Boys' Human Capital
Outcomes? Evidence from a Natural Experiment. \textit{Journal of Public Economics}
92(10--11): 2083--2105.}
{Compares blind (external) and non-blind (teacher) exam scores for the same students to
recover grading bias; finds bias against boys in this setting. Foundational
quasi-experimental template for detecting evaluator bias.}{RF}
% \note{}

\entry{Lavy, V., and Sand, E. (2018). On the Origins of Gender Gaps in Human Capital:
Short- and Long-Term Consequences of Teachers' Biases. \textit{Journal of Public
Economics} 167: 263--279.}
{Early-grade teacher grading bias has lasting effects on later course choices and
achievement, tracing a long-run channel from biased evaluation to specialization.
Directly relevant to STEM/non-STEM sorting.}{RF}
% \note{}


% =====================================================================
\section{D. How do students and parents respond to recommendations?}
\secintro{Compliance, the effect of making recommendations binding vs.\ advisory, and
SES heterogeneity --- mostly reform-based natural experiments.}

\entry{Bach, M. (2023). Heterogeneous Responses to School Track Choice: Evidence from
the Repeal of Binding Track Recommendations. \textit{Economics of Education Review}.}
{Exploits German states' repeal of binding recommendations (binding $\to$ signal). Free
choice does not widen the SES gap in academic-track enrolment, because previously-eligible
low-SES students also move up; a key mechanism is preferences over the intermediate
track. The cleanest design for separating the binding/authority arm from the
informational arm.}{RF}
\note{Reform goes binding $\to$ non-binding, so the recommendation becomes a pure signal --
the regime Flanders is already in for the soft \emph{advies}. Useful contrast: Bach gets
binding-vs-signal \emph{across a reform in time}; Belgium gives it \emph{within one system}
via two coexisting instruments (non-binding advies + binding attest). NOT ``Germany became
like Belgium'' -- Belgium has both arms at once. Mechanism (not pushing anyone down): two
offsetting flows at different margins. Ineligible high-SES move UP into academic (aspiration
-- would widen gap). But eligible low-SES \emph{also} move up, offsetting it, because free
choice floods the INTERMEDIATE track with low achievers from the basic track; eligible
low-SES anticipate worse middle-track peers and flee UPWARD to academic. So a worsening
middle \emph{repels} those who can go higher -- a GE/peer-composition effect operating
through track composition, not own-information. Direct support for keeping the peer fixed
point $\Theta_{sj}$ live; warning that my CFs are currently partial-equilibrium. Pairs with
OPS (parents too aspirational, teacher corrective) and de Ree (teacher too conservative,
parent corrective) -- opposite conclusions because none observe beliefs/match quality.}


\entry{Piopiunik, M. (2014). The Effects of Early Tracking on Student Performance:
Evidence from a School Reform in Bavaria. \textit{Economics of Education Review} 42:
12--33.}
{Reform-based evidence on how track placement (in a recommendation-driven system)
affects performance. Background on the stakes of the placement recommendations steer.}{RF}
\note{revise memoire, not that useful I think}


% =====================================================================
\section{F. Belief formation and information frictions in educational choice}
\secintro{The framework for how a signal (grade, recommendation, disclosed information)
shifts beliefs and choices --- awareness, ambiguity, learning about comparative advantage,
information experiments. Some are structural; flagged.}

\entry{Giustinelli, P., and Pavoni, N. (2017). The Evolution of Awareness and Belief
Ambiguity in the Process of High School Track Choice. \textit{Review of Economic
Dynamics} 25: 93--120.}
{Repeated survey of Italian 8th-graders: students are only partially aware of available
tracks and hold ambiguous beliefs about completing each, with disadvantaged students
acquiring information more slowly. Direct evidence on the belief object a recommendation
acts on, and on SES heterogeneity in updating.}{RF / survey}
\note{Giustinelli–Pavoni directly elicit students' awareness and ambiguous subjective beliefs about track success over the months before a non-binding, family-choice tracking decision, and show these beliefs are SES-graded, evolve through incomplete and selective learning, and predict enrollment — establishing the demand-side object (a learning student with measurable, ambiguous beliefs) that the teacher-recommendation literature assumes away; we connect that learning student to the teacher's recommendation as the signal that updates her beliefs, which they measure but never link to the recommendation itself.}

\entry{Giustinelli, P. (2016). Group Decision Making with Uncertain Outcomes: Unpacking
Child--Parent Choice of the High School Track. \textit{International Economic Review}.}
{Models and estimates the joint child--parent track decision using elicited subjective
expectations, quantifying how information/sensitization campaigns would shift the track
distribution. Relevant if parental influence matters.}{STR / RF}
\note{bargaining model over track choice}

\entry{Wiswall, M., and Zafar, B. (2015). Determinants of College Major Choice:
Identification Using an Information Experiment. \textit{Review of Economic Studies}
82(2): 791--824.}
{Provides students true population information on major-specific outcomes, generating
experimental belief variation; students revise logically and choices respond, though
tastes dominate. The template for treating an external signal as a belief shifter about
returns.}{EXP}
\note{reference for information problems... concludes that most of the preference is realized before H.E, which justifies our approach to understanding how preferences are shaped in the early stages when there are even less information + more ambiguity over taste, returns etc.}

\entry{Wiswall, M., and Zafar, B. (2015). How Do College Students Respond to Public
Information about Earnings? \textit{Journal of Human Capital} 9(2).}
{Companion: most students are non-Bayesian updaters, and assuming Bayesian processing
substantially overstates the welfare gains from information. Important caution for any
moment-matching/Bayesian-updating assumption about how students absorb advice.}{EXP}
% \note{}

\entry{Hastings, J., Neilson, C., and Zimmerman, S. (2015). The Effects of Earnings
Disclosure on College Enrollment Decisions. NBER WP; and Hastings, Neilson, and Zimmerman
(2013) on returns to majors.}
{Field experiment disclosing program-level earnings: muted average enrolment effects, but
low-SES applicants shift toward higher-net-return degrees. Real-world analog of an
informed party moving disadvantaged students' choices.}{EXP / RF}
\note{all these papers imply information at the earning level and are applied to H.E setting}


\entry{Aucejo, E., and James, J. (2021). The Path to College Education: The Role of Math
and Verbal Skills. \textit{Journal of Political Economy} 129(10): 2905--2946.}
{Shows relative (math vs.\ verbal) comparative advantage shapes the path into college and
field, beyond absolute skill. Useful microfoundation for STEM/non-STEM sorting on
comparative advantage --- the dimension your subtrack recommendation encodes.}{STR / RF}
\note{Aucejo–James provide the dynamic latent-factor methodology and the structural evidence that school quality is a skill-producing input with positive skill-environment complementarity — the engine behind our skill-accumulation block and the "track builds skill on top of endowment" mechanism; we adopt their measurement-and-production architecture but embed it in a model of endogenous, signal-driven track choice, so that the environment is selected rather than given, and we can evaluate whether sorting into skill-producing tracks is well-targeted rather than only how skill is produced. Useful for the production function\dots}


\entry{Stinebrickner, T., and Stinebrickner, R. (2012; 2014). Learning about Academic
Ability and the College Dropout Decision / Academic Performance and College Dropout.
\textit{Journal of Labor Economics} 30(4) and 32(3): 601--644.}
{Foundational learning models using longitudinal expectations data: students update
Bayesianly about ability from grades and revise enrolment. The canonical ``learning about
own ability from signals'' framework the recommendation extends.}{STR}
\note{reference on learning and unobserved heterogeneity.}

\entry{Arcidiacono, P., Aucejo, E., Maurel, A., and Ransom, T. (2025). College Attrition
and the Dynamics of Information Revelation. \textit{Journal of Political Economy} 133(1):
53--110.}
{State-of-the-art dynamic model in which students learn about match quality from signals
while choosing majors and attrition. The structural backbone many recommendation models
build on (signals arrive exogenously here --- the gap a strategic-school model fills).}{STR}
\note{reference on learning and unobserved heterogeneity. Main Machinery to get identification in the learning parameters}





\entry{Larroucau, Ríos; DYNAMIC COLLEGE ADMISSIONS, 2025}
{We study the relevance of incorporating dynamic incentives and eliciting private information about students’ preferences to improve their welfare and downstream outcomes in centralized assignment mechanisms. Using administrative data and two nationwide surveys, we identify two behavioral channels that largely explain students’ dynamic decisions: (i) initial mismatches and (ii) learning. Based on these facts, we build and estimate a structural model of students’ college progression in the presence of a centralized admission system, allowing students to learn about their preferences and abilities over time and reapply to the system. We use the estimated model to analyze the impact of changing the information environment and the assignment mechanism and reapplication rules on the efficiency of the system. Our counterfactual results show that policies that front-load learning can help to reduce switches, while providing score bonuses that elicit information on students’ cardinal preferences and leverage dynamic incentives can significantly increase retention and students’ overall welfare.}{STR}
\note{reference on learning and unobserved heterogeneity. Close to AAMR, but with an added twist that they allow learning to depend heterogeneously onSES, gender etc, nice to take as reference}



\entry{Bunting, Diegert, Maurel; DYNAMIC COLLEGE ADMISSIONS, 2025}
{We study the relevance of incorporating dynamic incentives and eliciting private information about students’ preferences to improve their welfare and downstream outcomes in centralized assignment mechanisms. Using administrative data and two nationwide surveys, we identify two behavioral channels that largely explain students’ dynamic decisions: (i) initial mismatches and (ii) learning. Based on these facts, we build and estimate a structural model of students’ college progression in the presence of a centralized admission system, allowing students to learn about their preferences and abilities over time and reapply to the system. We use the estimated model to analyze the impact of changing the information environment and the assignment mechanism and reapplication rules on the efficiency of the system. Our counterfactual results show that policies that front-load learning can help to reduce switches, while providing score bonuses that elicit information on students’ cardinal preferences and leverage dynamic incentives can significantly increase retention and students’ overall welfare.}{STR}
\note{reference on learning and unobserved heterogeneity. Close to AAMR, but with an added twist that they allow learning to depend heterogeneously onSES, gender etc, nice to take as reference}



% =====================================================================
\section{H. Tracking: foundational context}
\secintro{The broader tracking literature situating why placement matters and what is at
stake in the recommendation.}

\entry{Card, D., and Giuliano, L. (2016). Can Tracking Raise the Test Scores of
High-Ability Minority Students? \textit{American Economic Review} 106(10): 2783--2816.}
{Ability tracking/placement can raise outcomes for high-ability minority students,
highlighting that who is placed where (and how it is decided) has real consequences.}{RF}
% \note{}

\entry{Dustmann, C., Puhani, P. A., and Sch\"onberg, U. (2017). The Long-Term Effects of
Early Track Choice. \textit{The Economic Journal} 127.}
{Estimates long-run labour-market and attainment effects of early between-school tracking,
with attention to track mobility. Motivates the stakes and irreversibility of the age-14
decision.}{RF}
% \note{}

\entry{Hanushek, E. A., and Woessmann, L. (2006). Does Educational Tracking Affect
Performance and Inequality? Differences-in-Differences Evidence Across Countries.
\textit{The Economic Journal} 116(510): C63--C76.}
{Cross-country diff-in-diff: early tracking tends to increase inequality without clear
average gains. Canonical macro evidence on tracking and inequality.}{RF}
% \note{}

\entry{De Groote, O. (2025). Dynamic Effort Choice in High School: Costs and Benefits of
an Academic Track. \textit{Journal of Labor Economics} 43(2): 467--502.}
{Structural treatment of the same Belgian/LOSO setting (effort, track, performance
dynamics). The natural in-house anchor for data and institutional detail.}{STR}
% \note{}

\entry{Betts, J. R. (2011). The Economics of Tracking in Education. In \textit{Handbook
of the Economics of Education}, vol.\ 3: 341--381.}
{Survey of tracking designs, identification challenges, and findings. Good orienting
reference for the tracking literature as a whole.}{REV}
% \note{}

\textbf{Add Fu Mehta\dots}


\section{H. Theory papers}

Here I would add cheap talk and bayesian persuasion to our setting.



\section{I. On the importance of STEM specialization in high-school}
The idea behind this is to understand tracking as not only in the usual vertical dimension of skill but also along different specializations. More and more the literature has been producing evidence at this STEM-nonSTEM dimension. Here is abstracting how these decisions are done, under uncertainty, under information settings etc but simply is about how consequential these decisions can be. To cite two recent examples at the frontier of education:

\entry{Ainsworth, Dehejia, Munteanu, Pop-Eleches, Urquiola. Malleable minds: The effects of STEM vs. humanities-focused curricula (2026, WP)}
{We examine the impacts of assignment to STEM vs. humanities-focused curricula in Romania’s high school system. We apply a regression discontinuity design to administrative and survey data to estimate effects on educational pathways, desired careers, and noncognitive outcomes. An overarching theme of our findings is the malleability of students to what they study. Assignment to STEM increases STEM college enrollment and technology or engineering career intentions by 25 pp. Exploring mechanisms, we find that STEM assignment changes students’ self-perceived academic abilities and their preferences over academic subjects and job tasks. STEM assignment is risky for low-achieving students, reducing their chances of passing a high school exit exam and enrolling in college. A final finding is that STEM makes boys more conservative, while shifting some of girls’ views to the left. Our results identify a strategy for promoting STEM higher education and careers, but also highlight potential tradeoffs.}{REV}
\note{In Rumania, very similar setting as in Belgium where tracks are implemented with some rigidities to transfer and where higher education is free and universal and not mediated by entry exams etc. MAIN FINDING: Our overarching finding is that high school students are malleable to what they study. We first show that high school curricular assignment has large effects on college and career choices in Romania, consistent with studies in other settings. We then establish four novel results. First, curricular assignment also affects what students enjoy and think they are good at, and these effects appear to drive the college and career impacts. Second, focusing on STEM rather than humanities causes students to report greater wellbeing in high school—despite spending more time on schoolwork—although this boost fades after graduation,, THIRD, doing stem more likely to fail... \dots In sum, we provide extensive evidence that curricular assignment affects students’ interests and beliefs about their abilities—and that these are key factors in their college and career choices\dots our results match the existing literature (Section 2) in showing that high school tracks and coursework affect educational pathways and labor market outcomes.37 Further, the finding of heterogeneity for effects on career choices—but not college choices—is consistent with evidence that career choices are subject to additional considerations, such as tastes for work hours and conditions, which may vary across student types (Patnaik, Wiswall, and Zafar 2020).}


\entry{John Eric Humphries Juanna Schrøter Joensen Gregory F. Veramendi. Complementarities in High School  and College Investments (2026, WP)}
{This paper examines how high school specialization shapes college investment decisions and their subsequent returns through dynamic complementarities. Using Swedish administrative data, we estimate a dynamic Roy model that accounts for selection on multidimensional skills, family background, prior investments, and unobserved heterogeneity. We identify the model using rich skill measures and quasi-experimental variation in program popularity. For marginal students, STEM specialization in high school increases wages by 9\%, with more than half this return attributed to dynamic complementarities that enhance the productivity of subsequent college investments. Consequently, we find that counterfactual policies encouraging high school STEM specialization generate twice the returns of equivalent college-level interventions. These findings demonstrate how the timing of specialized human capital investments matters during adolescence, with important implications for education policies that encourage or restrict specialization.}{REV}
\note{Fairly similar sub-track choice, fairly similar information and identification of latent factors as us. Structurally estimate the dynamics of the production function, like Attanasio et al (2020): .This is simply an extension of Heckman et al 2018 semi-structural model. We document large but heterogeneous returns from specializing in STEM in high school, and show that around half of the return comes from complementarities between high school and college investments.The model includes both specialization decisions (i.e., track and major) and attainment in high school and college. Using the model, we document rich sorting on multidimensional skills into high school track, followed by sorting on both skills and high school track into college majors. Second, we use the model to estimate and decompose treatment effects from specializing in specific high school tracks in high school. We find that the returns to the academic STEM track are high on average, but do not benefit  39 everyone. We then decompose the treatment effects into a direct effect, changes in college enrollment and major decisions, and complementarities between high school and college choices. For the academic STEM track, complementarities account for around 60\% of the wage gains, with changes in college and direct effects each accounting for around 20\%. Third, we use the model to evaluate two counterfactual policies designed to promote marginal STEM enrollment either in high school or when applying to college. We find that both colleges increase the number of STEM enrollees and graduates in college, but that the high school policy creates larger wage gains and benefits a greater share of those affected.  Our findings highlight a fundamental trade-off in education system design: early specialization can enhance returns through dynamic complementarities for students who maintain a consistent specialization path, but may change options and reduce returns for students who switch fields. This suggests that policies encouraging early specialization should be accompanied by flexibility for students to adjust their education paths as they discover their comparative advantages. More broadly, our findings demonstrate that understanding the interplay between skills, the timing of specialized investments, and the constraints in education systems is crucial for developing human capital policies that improve both individual outcomes and economic efficiency.}

CES etc... the question is do we need such thing?

\entry{Gordon B. Dahl, Dan-Olof Rooth, Anders Stenberg. High School Majors and Future Earnings (2021?)}
{We study how high school majors affect adult earnings using a regression discontinuity design. In Sweden, students are admitted to majors in tenth grade based on their
preference rankings and ninth grade GPA. We find Engineering, Natural Science, and Business
majors yield higher earnings than Social Science and Humanities, with major-specific returns
also varying based on next-best alternatives. There is either a zero or a negative return to
completing an academic program for students with a second-best non-academic major. Most
of the differences in adult earnings can be attributed to differences in occupation, and to a
lesser extent, college major}{REV}
\note{While this paper makes important progress on estimating long-term payoffs to high school
majors, several questions remain unanswered. The parameters we estimate are ex-post payoffs
to majors. An interesting question for future research is whether these ex-post payoffs line
up with ex-ante predicted payoffs.27 If they do, it suggests that students understand the
monetary tradeoffs associated with different majors, and that some students are willing to
trade off higher earnings for non-pecuniary returns. However, it is also possible that at
age 16, students do not yet know what occupation will be the best fit for them and they
may not be knowledgeable about earnings differences across fields. In future work, it would
be interesting to explore the factors influencing an individual’s major choice, including the
impact of parents, friends, and teachers. The parameters we estimate are also for compliers
on the margin of gaining entry into a major. For these marginal individuals, the effects can
be as large in absolute value as the returns to 2 years of additional schooling. It would be
interesting to know if similar patterns hold for other individuals.
}



% =====================================================================
\vspace{6pt}
\hrule
\vspace{6pt}
\noindent\textcolor{grey}{\small \textbf{Notes on use.} Most on-theme empirical anchors:
schools-know-more --- Jackson (2018), Kinsler--Pavan (2016), Dizon-Ross (2019), Boneva--Rauh
(2018); recommendations as causal --- Papageorge--Gershenson--Kang (2020), Diamond--Persson
(2016), Mulhern (2022); the class-rank channel for $Z$ --- Murphy--Weinhardt (2020),
Elsner--Isphording (2017); bias/discrimination in how schools decide --- Carlana (2019) and
the blind-vs-non-blind designs (Botelho et al.\ 2015; Lavy 2008; Lavy--Sand 2018); how
students respond --- Bach (2023) and the Dutch cutoff designs; interpretive toolkit ---
Bohren et al.\ (2025), Altonji--Pierret (2001), Arcidiacono--Bayer--Hizmo (2010). Items
tagged STR/THY are for panorama and framing rather than reduced-form templates. Confirm any
venue marked \emph{(verify)} before citing.}

\end{document}s